% ============================================================================
% CV template: Minimal violet (two pages at most, from Work Finder and Atelier)
% Layout only. Every name, contact, headline and bullet lives in the content
% file, which tailor-tex generates from a CvDocument for each build; every
% text value there is escaped by tailor-tex, never written by a model.
%
% Compile with XeLaTeX, twice:
%   xelatex "\def\ContentFile{content}\def\FontDir{fonts/}% ============================================================================
% CV template: Minimal violet (two pages at most, from Work Finder and Atelier)
% Layout only. Every name, contact, headline and bullet lives in the content
% file, which tailor-tex generates from a CvDocument for each build; every
% text value there is escaped by tailor-tex, never written by a model.
%
% Compile with XeLaTeX, twice:
%   xelatex "\def\ContentFile{content}\def\FontDir{fonts/}% ============================================================================
% CV template: Minimal violet (two pages at most, from Work Finder and Atelier)
% Layout only. Every name, contact, headline and bullet lives in the content
% file, which tailor-tex generates from a CvDocument for each build; every
% text value there is escaped by tailor-tex, never written by a model.
%
% Compile with XeLaTeX, twice:
%   xelatex "\def\ContentFile{content}\def\FontDir{fonts/}% ============================================================================
% CV template: Minimal violet (two pages at most, from Work Finder and Atelier)
% Layout only. Every name, contact, headline and bullet lives in the content
% file, which tailor-tex generates from a CvDocument for each build; every
% text value there is escaped by tailor-tex, never written by a model.
%
% Compile with XeLaTeX, twice:
%   xelatex "\def\ContentFile{content}\def\FontDir{fonts/}\input{cv}"
% \ContentFile defaults to "content", \FontDir to "fonts/" (fonts/ibm-plex-sans
% of the repo). tailor-tex does all of this in a temporary folder per build.
% An empty job location drops its middle dot instead of
% leaving it dangling, and a summary with an empty second paragraph drops
% the gap before it.
%
% Font: IBM Plex Sans (OFL, fonts/ibm-plex-sans/OFL.txt). Type scale (pt):
% name 29, headline 12.5, job title 10.6, body 9.5 on 13, metadata 8.8,
% section labels 9.2 tracked. Accent violet #6D28D9.
%
% ATS rules this layout keeps (verified with poppler and a drawing-order
% parser): single text column, real text only, no tables, icons, header or
% footer; no ligatures and no hyphenation; a font that maps every glyph to
% plain ASCII (hyphens and parentheses included); tracking at or below
% 0.06 em; skills one category per line; inline items joined by middle dots,
% never by wide gaps; nothing right-aligned alone where a parser would read
% it as a separate column.
% ============================================================================
\documentclass[a4paper,10pt]{article}

\providecommand{\ContentFile}{content}
\providecommand{\FontDir}{fonts/}

% ----------------------------------------------------------------------------
% Page and packages
% ----------------------------------------------------------------------------
\usepackage[a4paper,left=15mm,right=15mm,top=13mm,bottom=14mm]{geometry}
\usepackage{fontspec}
\usepackage[english]{babel}
\usepackage[protrusion=true]{microtype}
\usepackage{enumitem}
\usepackage{titlesec}
\usepackage{needspace}
\usepackage{etoolbox}
\usepackage{xcolor}
\usepackage[hidelinks,unicode]{hyperref}

% ----------------------------------------------------------------------------
% Fonts: plain letters only, so "fi", "fl" and "ff" stay searchable. IBM Plex
% Sans, bold text in SemiBold; it keeps I, l and 1 distinct, so OCR-based
% checkers read it well.
% ----------------------------------------------------------------------------
\def\CvRegular{IBMPlexSans-Regular.ttf}
\def\CvBold{IBMPlexSans-SemiBold.ttf}
\def\CvHeavy{IBMPlexSans-Bold.ttf}
\def\CvNameTrack{-0.5}
\defaultfontfeatures{Ligatures={TeX,NoCommon},Kerning=On}
\setmainfont{\CvRegular}[Path=\FontDir,BoldFont=\CvBold]
\newfontfamily\namefont{\CvHeavy}[Path=\FontDir,LetterSpace=\CvNameTrack]
% Tracking stays below 0.1 em, so parsers still read every label as one word.
\newfontfamily\labelfont{\CvBold}[Path=\FontDir,LetterSpace=6]

% No hyphenation: every keyword stays whole for ATS parsers.
\hyphenpenalty=10000
\exhyphenpenalty=10000
\setlength{\emergencystretch}{3em}

% ----------------------------------------------------------------------------
% Colors: one violet accent, a deep violet-ink for highlights, neutral text
% ----------------------------------------------------------------------------
\providecommand{\AccentHex}{6D28D9}
\definecolor{Accent}{HTML}{\AccentHex}
\colorlet{Highlight}{Accent!32!black}
\colorlet{Rule}{Accent!18!white}
\colorlet{Divider}{Accent!45!white}
\definecolor{Ink}{HTML}{14171C}
\definecolor{Body}{HTML}{2B3038}
\definecolor{Muted}{HTML}{646B76}
\setlength{\parindent}{0pt}
\setlength{\parskip}{0pt}
\raggedbottom
% Never split a bullet so that one of its lines sits alone on a page.
\clubpenalty=10000
\widowpenalty=10000
\displaywidowpenalty=10000

% No header or footer: running text there lands in the middle of a job's
% bullets when ATS parsers extract the page.
\pagestyle{empty}

% ----------------------------------------------------------------------------
% Header in three rows: the name with the title at the right edge on its
% baseline; then location (with time zone and work format) and the phone,
% with the core stack at the right edge; then every link on one line. Each of
% location, phone and stack may be empty; with all three empty the row is gone.
% Links use the accent so they read as clickable.
%   \cvheader{name}{location}{phone as shown}{phone digits}{headline}{stack}
%   \cvlinks{\cvlink{url}{text}\linksep\cvlink{url}{text}...}
% The links never wrap: a row too wide for the line first gets narrower gaps
% (7pt down to 4pt), then a smaller size (9.2pt down to 7.6pt). Only past
% that floor does it wrap, rather than turn unreadable.
% ----------------------------------------------------------------------------
\newlength{\cvlinkgap}
\setlength{\cvlinkgap}{7pt}
\newlength{\cvlinksize}
\newlength{\cvlinkswidth}
\newif\ifcvlinksshrink
\newcommand{\linksep}{\hspace{\cvlinkgap}{\color{Divider}·}\hspace{\cvlinkgap}}
\newcommand{\cvlink}[2]{\href{#1}{\color{Accent}#2}}
\newcommand{\cvheader}[6]{%
  \hypersetup{pdfauthor={#1}}%
  \noindent{\namefont\fontsize{29}{33}\selectfont\color{Ink}#1}\hfill
  {\fontsize{12.5}{16}\selectfont\bfseries\color{Accent}#5}\par
  \ifboolexpr{test {\ifblank{#2}} and test {\ifblank{#3}} and test {\ifblank{#6}}}{\vspace{4pt}}{%
    \vspace{2pt}%
    {\noindent\fontsize{9.5}{13}\selectfont
      \ifblank{#2}{}{{\color{Ink}#2}}%
      \ifboolexpr{test {\ifblank{#2}} or test {\ifblank{#3}}}{}{\linksep}%
      \ifblank{#3}{}{{\color{Ink}\href{tel:#4}{#3}}}%
      \hfill{\color{Muted}#6}\par}%
    \vspace{5pt}}}
% Width of the row at the current gap and size, into \cvlinkswidth.
\newcommand{\cvlinksmeasure}[1]{%
  \settowidth{\cvlinkswidth}{\fontsize{\cvlinksize}{12}\selectfont#1}}
% One step of the fit by recursion, not \loop: font loading at a new size
% may run a \loop of its own and clobber ours.
\newcommand{\cvlinksfit}[1]{%
  \cvlinksshrinkfalse
  \ifdim\cvlinkswidth>\linewidth
    \ifdim\cvlinkgap>4pt
      \addtolength{\cvlinkgap}{-0.5pt}\cvlinksshrinktrue
    \else\ifdim\cvlinksize>7.6pt
      \addtolength{\cvlinksize}{-0.2pt}\cvlinksshrinktrue
    \fi\fi
  \fi
  \ifcvlinksshrink
    \def\cvlinksnext{\cvlinksmeasure{#1}\cvlinksfit{#1}}%
  \else
    \let\cvlinksnext\relax
  \fi
  \cvlinksnext}
\newcommand{\cvlinks}[1]{{%
  \setlength{\cvlinkgap}{7pt}%
  \setlength{\cvlinksize}{9.2pt}%
  \cvlinksmeasure{#1}%
  \cvlinksfit{#1}%
  \fontsize{\cvlinksize}{12}\selectfont#1\par}\vspace{2pt}}

% ----------------------------------------------------------------------------
% Sections and blocks
% ----------------------------------------------------------------------------
% Section: tracked uppercase label followed by a hairline to the margin.
\newcommand{\sectionlabel}[1]{\MakeUppercase{#1}\hspace{10pt}{\color{Rule}\leaders\vrule height 3.5pt depth -2.85pt\hfill}\kern0pt}
\titleformat{\section}{\labelfont\fontsize{9.2}{12}\selectfont\color{Accent}}{}{0pt}{\sectionlabel}
\titlespacing*{\section}{0pt}{10pt}{5.5pt}
\newcommand{\cvsection}[1]{\Needspace{4\baselineskip}\section*{#1}}

\newcommand{\hl}[1]{{\bfseries\color{Highlight}#1}}

\newcommand{\cvsummary}[2]{#1\par\ifblank{#2}{}{\vspace{4.5pt}#2\par}}

% Skills: bold label and the list on the same line, one line per category.
\newenvironment{cvskills}{\par}{\par}
\newcommand{\cvskill}[2]{%
  \par\hangindent=10pt\hangafter=1\noindent
  {\bfseries\color{Ink}#1:}\hspace{4pt}#2\par\vspace{2.6pt}}

% Job: title | company on the left, dates and location on the right,
% a one-line company description underneath; never stranded at a page end.
%   \cvjob{title}{company}{dates}{location}{company description}
\newcommand{\cvjob}[5]{%
  \par\addvspace{8pt}\Needspace{4\baselineskip}%
  {\fontsize{10.6}{14}\selectfont{\bfseries\color{Ink}#1}\hspace{7pt}{\color{Divider}|}\hspace{7pt}{\bfseries\color{Accent}#2}\hfill{\fontsize{8.8}{14}\selectfont\color{Muted}#3\ifblank{#4}{}{\hspace{5pt}·\hspace{5pt}#4}}}\par
  {\fontsize{8.8}{12.4}\selectfont\color{Muted}#5}\par\nopagebreak
  \vspace{3.5pt}}

\newenvironment{cvbullets}{%
  \begin{itemize}[leftmargin=12.5pt,labelsep=6.5pt,itemsep=2pt,topsep=0pt,parsep=0pt,partopsep=0pt,
    label={\color{Accent}\fontsize{8.5}{8.5}\selectfont\textbullet}]}
  {\end{itemize}}

% Education: one line like a job's, degree | school, the dates (if any) at
% the right edge.
\newcommand{\cvedu}[3]{%
  {\fontsize{10.6}{14}\selectfont{\bfseries\color{Ink}#1}\hspace{7pt}{\color{Divider}|}\hspace{7pt}{\bfseries\color{Accent}#2}\hfill{\fontsize{8.8}{14}\selectfont\color{Muted}#3}}\par}

% Languages share one line, joined by a middle dot. A wide gap instead would
% make parsers read the second language as a separate column.
\newcommand{\cvlangsep}{}
\newcommand{\cvlang}[2]{\cvlangsep{\bfseries\color{Ink}#1:}\hspace{4pt}#2%
  \gdef\cvlangsep{\hspace{6pt}{\color{Muted}·}\hspace{6pt}}}

% ============================================================================
\begin{document}
\fontsize{9.5}{13}\selectfont
\color{Body}
\raggedright
\frenchspacing

\input{\ContentFile}

\end{document}
"
% \ContentFile defaults to "content", \FontDir to "fonts/" (fonts/ibm-plex-sans
% of the repo). tailor-tex does all of this in a temporary folder per build.
% An empty job location drops its middle dot instead of
% leaving it dangling, and a summary with an empty second paragraph drops
% the gap before it.
%
% Font: IBM Plex Sans (OFL, fonts/ibm-plex-sans/OFL.txt). Type scale (pt):
% name 29, headline 12.5, job title 10.6, body 9.5 on 13, metadata 8.8,
% section labels 9.2 tracked. Accent violet #6D28D9.
%
% ATS rules this layout keeps (verified with poppler and a drawing-order
% parser): single text column, real text only, no tables, icons, header or
% footer; no ligatures and no hyphenation; a font that maps every glyph to
% plain ASCII (hyphens and parentheses included); tracking at or below
% 0.06 em; skills one category per line; inline items joined by middle dots,
% never by wide gaps; nothing right-aligned alone where a parser would read
% it as a separate column.
% ============================================================================
\documentclass[a4paper,10pt]{article}

\providecommand{\ContentFile}{content}
\providecommand{\FontDir}{fonts/}

% ----------------------------------------------------------------------------
% Page and packages
% ----------------------------------------------------------------------------
\usepackage[a4paper,left=15mm,right=15mm,top=13mm,bottom=14mm]{geometry}
\usepackage{fontspec}
\usepackage[english]{babel}
\usepackage[protrusion=true]{microtype}
\usepackage{enumitem}
\usepackage{titlesec}
\usepackage{needspace}
\usepackage{etoolbox}
\usepackage{xcolor}
\usepackage[hidelinks,unicode]{hyperref}

% ----------------------------------------------------------------------------
% Fonts: plain letters only, so "fi", "fl" and "ff" stay searchable. IBM Plex
% Sans, bold text in SemiBold; it keeps I, l and 1 distinct, so OCR-based
% checkers read it well.
% ----------------------------------------------------------------------------
\def\CvRegular{IBMPlexSans-Regular.ttf}
\def\CvBold{IBMPlexSans-SemiBold.ttf}
\def\CvHeavy{IBMPlexSans-Bold.ttf}
\def\CvNameTrack{-0.5}
\defaultfontfeatures{Ligatures={TeX,NoCommon},Kerning=On}
\setmainfont{\CvRegular}[Path=\FontDir,BoldFont=\CvBold]
\newfontfamily\namefont{\CvHeavy}[Path=\FontDir,LetterSpace=\CvNameTrack]
% Tracking stays below 0.1 em, so parsers still read every label as one word.
\newfontfamily\labelfont{\CvBold}[Path=\FontDir,LetterSpace=6]

% No hyphenation: every keyword stays whole for ATS parsers.
\hyphenpenalty=10000
\exhyphenpenalty=10000
\setlength{\emergencystretch}{3em}

% ----------------------------------------------------------------------------
% Colors: one violet accent, a deep violet-ink for highlights, neutral text
% ----------------------------------------------------------------------------
\providecommand{\AccentHex}{6D28D9}
\definecolor{Accent}{HTML}{\AccentHex}
\colorlet{Highlight}{Accent!32!black}
\colorlet{Rule}{Accent!18!white}
\colorlet{Divider}{Accent!45!white}
\definecolor{Ink}{HTML}{14171C}
\definecolor{Body}{HTML}{2B3038}
\definecolor{Muted}{HTML}{646B76}
\setlength{\parindent}{0pt}
\setlength{\parskip}{0pt}
\raggedbottom
% Never split a bullet so that one of its lines sits alone on a page.
\clubpenalty=10000
\widowpenalty=10000
\displaywidowpenalty=10000

% No header or footer: running text there lands in the middle of a job's
% bullets when ATS parsers extract the page.
\pagestyle{empty}

% ----------------------------------------------------------------------------
% Header in three rows: the name with the title at the right edge on its
% baseline; then location (with time zone and work format) and the phone,
% with the core stack at the right edge; then every link on one line. Each of
% location, phone and stack may be empty; with all three empty the row is gone.
% Links use the accent so they read as clickable.
%   \cvheader{name}{location}{phone as shown}{phone digits}{headline}{stack}
%   \cvlinks{\cvlink{url}{text}\linksep\cvlink{url}{text}...}
% The links never wrap: a row too wide for the line first gets narrower gaps
% (7pt down to 4pt), then a smaller size (9.2pt down to 7.6pt). Only past
% that floor does it wrap, rather than turn unreadable.
% ----------------------------------------------------------------------------
\newlength{\cvlinkgap}
\setlength{\cvlinkgap}{7pt}
\newlength{\cvlinksize}
\newlength{\cvlinkswidth}
\newif\ifcvlinksshrink
\newcommand{\linksep}{\hspace{\cvlinkgap}{\color{Divider}·}\hspace{\cvlinkgap}}
\newcommand{\cvlink}[2]{\href{#1}{\color{Accent}#2}}
\newcommand{\cvheader}[6]{%
  \hypersetup{pdfauthor={#1}}%
  \noindent{\namefont\fontsize{29}{33}\selectfont\color{Ink}#1}\hfill
  {\fontsize{12.5}{16}\selectfont\bfseries\color{Accent}#5}\par
  \ifboolexpr{test {\ifblank{#2}} and test {\ifblank{#3}} and test {\ifblank{#6}}}{\vspace{4pt}}{%
    \vspace{2pt}%
    {\noindent\fontsize{9.5}{13}\selectfont
      \ifblank{#2}{}{{\color{Ink}#2}}%
      \ifboolexpr{test {\ifblank{#2}} or test {\ifblank{#3}}}{}{\linksep}%
      \ifblank{#3}{}{{\color{Ink}\href{tel:#4}{#3}}}%
      \hfill{\color{Muted}#6}\par}%
    \vspace{5pt}}}
% Width of the row at the current gap and size, into \cvlinkswidth.
\newcommand{\cvlinksmeasure}[1]{%
  \settowidth{\cvlinkswidth}{\fontsize{\cvlinksize}{12}\selectfont#1}}
% One step of the fit by recursion, not \loop: font loading at a new size
% may run a \loop of its own and clobber ours.
\newcommand{\cvlinksfit}[1]{%
  \cvlinksshrinkfalse
  \ifdim\cvlinkswidth>\linewidth
    \ifdim\cvlinkgap>4pt
      \addtolength{\cvlinkgap}{-0.5pt}\cvlinksshrinktrue
    \else\ifdim\cvlinksize>7.6pt
      \addtolength{\cvlinksize}{-0.2pt}\cvlinksshrinktrue
    \fi\fi
  \fi
  \ifcvlinksshrink
    \def\cvlinksnext{\cvlinksmeasure{#1}\cvlinksfit{#1}}%
  \else
    \let\cvlinksnext\relax
  \fi
  \cvlinksnext}
\newcommand{\cvlinks}[1]{{%
  \setlength{\cvlinkgap}{7pt}%
  \setlength{\cvlinksize}{9.2pt}%
  \cvlinksmeasure{#1}%
  \cvlinksfit{#1}%
  \fontsize{\cvlinksize}{12}\selectfont#1\par}\vspace{2pt}}

% ----------------------------------------------------------------------------
% Sections and blocks
% ----------------------------------------------------------------------------
% Section: tracked uppercase label followed by a hairline to the margin.
\newcommand{\sectionlabel}[1]{\MakeUppercase{#1}\hspace{10pt}{\color{Rule}\leaders\vrule height 3.5pt depth -2.85pt\hfill}\kern0pt}
\titleformat{\section}{\labelfont\fontsize{9.2}{12}\selectfont\color{Accent}}{}{0pt}{\sectionlabel}
\titlespacing*{\section}{0pt}{10pt}{5.5pt}
\newcommand{\cvsection}[1]{\Needspace{4\baselineskip}\section*{#1}}

\newcommand{\hl}[1]{{\bfseries\color{Highlight}#1}}

\newcommand{\cvsummary}[2]{#1\par\ifblank{#2}{}{\vspace{4.5pt}#2\par}}

% Skills: bold label and the list on the same line, one line per category.
\newenvironment{cvskills}{\par}{\par}
\newcommand{\cvskill}[2]{%
  \par\hangindent=10pt\hangafter=1\noindent
  {\bfseries\color{Ink}#1:}\hspace{4pt}#2\par\vspace{2.6pt}}

% Job: title | company on the left, dates and location on the right,
% a one-line company description underneath; never stranded at a page end.
%   \cvjob{title}{company}{dates}{location}{company description}
\newcommand{\cvjob}[5]{%
  \par\addvspace{8pt}\Needspace{4\baselineskip}%
  {\fontsize{10.6}{14}\selectfont{\bfseries\color{Ink}#1}\hspace{7pt}{\color{Divider}|}\hspace{7pt}{\bfseries\color{Accent}#2}\hfill{\fontsize{8.8}{14}\selectfont\color{Muted}#3\ifblank{#4}{}{\hspace{5pt}·\hspace{5pt}#4}}}\par
  {\fontsize{8.8}{12.4}\selectfont\color{Muted}#5}\par\nopagebreak
  \vspace{3.5pt}}

\newenvironment{cvbullets}{%
  \begin{itemize}[leftmargin=12.5pt,labelsep=6.5pt,itemsep=2pt,topsep=0pt,parsep=0pt,partopsep=0pt,
    label={\color{Accent}\fontsize{8.5}{8.5}\selectfont\textbullet}]}
  {\end{itemize}}

% Education: one line like a job's, degree | school, the dates (if any) at
% the right edge.
\newcommand{\cvedu}[3]{%
  {\fontsize{10.6}{14}\selectfont{\bfseries\color{Ink}#1}\hspace{7pt}{\color{Divider}|}\hspace{7pt}{\bfseries\color{Accent}#2}\hfill{\fontsize{8.8}{14}\selectfont\color{Muted}#3}}\par}

% Languages share one line, joined by a middle dot. A wide gap instead would
% make parsers read the second language as a separate column.
\newcommand{\cvlangsep}{}
\newcommand{\cvlang}[2]{\cvlangsep{\bfseries\color{Ink}#1:}\hspace{4pt}#2%
  \gdef\cvlangsep{\hspace{6pt}{\color{Muted}·}\hspace{6pt}}}

% ============================================================================
\begin{document}
\fontsize{9.5}{13}\selectfont
\color{Body}
\raggedright
\frenchspacing

\input{\ContentFile}

\end{document}
"
% \ContentFile defaults to "content", \FontDir to "fonts/" (fonts/ibm-plex-sans
% of the repo). tailor-tex does all of this in a temporary folder per build.
% An empty job location drops its middle dot instead of
% leaving it dangling, and a summary with an empty second paragraph drops
% the gap before it.
%
% Font: IBM Plex Sans (OFL, fonts/ibm-plex-sans/OFL.txt). Type scale (pt):
% name 29, headline 12.5, job title 10.6, body 9.5 on 13, metadata 8.8,
% section labels 9.2 tracked. Accent violet #6D28D9.
%
% ATS rules this layout keeps (verified with poppler and a drawing-order
% parser): single text column, real text only, no tables, icons, header or
% footer; no ligatures and no hyphenation; a font that maps every glyph to
% plain ASCII (hyphens and parentheses included); tracking at or below
% 0.06 em; skills one category per line; inline items joined by middle dots,
% never by wide gaps; nothing right-aligned alone where a parser would read
% it as a separate column.
% ============================================================================
\documentclass[a4paper,10pt]{article}

\providecommand{\ContentFile}{content}
\providecommand{\FontDir}{fonts/}

% ----------------------------------------------------------------------------
% Page and packages
% ----------------------------------------------------------------------------
\usepackage[a4paper,left=15mm,right=15mm,top=13mm,bottom=14mm]{geometry}
\usepackage{fontspec}
\usepackage[english]{babel}
\usepackage[protrusion=true]{microtype}
\usepackage{enumitem}
\usepackage{titlesec}
\usepackage{needspace}
\usepackage{etoolbox}
\usepackage{xcolor}
\usepackage[hidelinks,unicode]{hyperref}

% ----------------------------------------------------------------------------
% Fonts: plain letters only, so "fi", "fl" and "ff" stay searchable. IBM Plex
% Sans, bold text in SemiBold; it keeps I, l and 1 distinct, so OCR-based
% checkers read it well.
% ----------------------------------------------------------------------------
\def\CvRegular{IBMPlexSans-Regular.ttf}
\def\CvBold{IBMPlexSans-SemiBold.ttf}
\def\CvHeavy{IBMPlexSans-Bold.ttf}
\def\CvNameTrack{-0.5}
\defaultfontfeatures{Ligatures={TeX,NoCommon},Kerning=On}
\setmainfont{\CvRegular}[Path=\FontDir,BoldFont=\CvBold]
\newfontfamily\namefont{\CvHeavy}[Path=\FontDir,LetterSpace=\CvNameTrack]
% Tracking stays below 0.1 em, so parsers still read every label as one word.
\newfontfamily\labelfont{\CvBold}[Path=\FontDir,LetterSpace=6]

% No hyphenation: every keyword stays whole for ATS parsers.
\hyphenpenalty=10000
\exhyphenpenalty=10000
\setlength{\emergencystretch}{3em}

% ----------------------------------------------------------------------------
% Colors: one violet accent, a deep violet-ink for highlights, neutral text
% ----------------------------------------------------------------------------
\providecommand{\AccentHex}{6D28D9}
\definecolor{Accent}{HTML}{\AccentHex}
\colorlet{Highlight}{Accent!32!black}
\colorlet{Rule}{Accent!18!white}
\colorlet{Divider}{Accent!45!white}
\definecolor{Ink}{HTML}{14171C}
\definecolor{Body}{HTML}{2B3038}
\definecolor{Muted}{HTML}{646B76}
\setlength{\parindent}{0pt}
\setlength{\parskip}{0pt}
\raggedbottom
% Never split a bullet so that one of its lines sits alone on a page.
\clubpenalty=10000
\widowpenalty=10000
\displaywidowpenalty=10000

% No header or footer: running text there lands in the middle of a job's
% bullets when ATS parsers extract the page.
\pagestyle{empty}

% ----------------------------------------------------------------------------
% Header in three rows: the name with the title at the right edge on its
% baseline; then location (with time zone and work format) and the phone,
% with the core stack at the right edge; then every link on one line. Each of
% location, phone and stack may be empty; with all three empty the row is gone.
% Links use the accent so they read as clickable.
%   \cvheader{name}{location}{phone as shown}{phone digits}{headline}{stack}
%   \cvlinks{\cvlink{url}{text}\linksep\cvlink{url}{text}...}
% The links never wrap: a row too wide for the line first gets narrower gaps
% (7pt down to 4pt), then a smaller size (9.2pt down to 7.6pt). Only past
% that floor does it wrap, rather than turn unreadable.
% ----------------------------------------------------------------------------
\newlength{\cvlinkgap}
\setlength{\cvlinkgap}{7pt}
\newlength{\cvlinksize}
\newlength{\cvlinkswidth}
\newif\ifcvlinksshrink
\newcommand{\linksep}{\hspace{\cvlinkgap}{\color{Divider}·}\hspace{\cvlinkgap}}
\newcommand{\cvlink}[2]{\href{#1}{\color{Accent}#2}}
\newcommand{\cvheader}[6]{%
  \hypersetup{pdfauthor={#1}}%
  \noindent{\namefont\fontsize{29}{33}\selectfont\color{Ink}#1}\hfill
  {\fontsize{12.5}{16}\selectfont\bfseries\color{Accent}#5}\par
  \ifboolexpr{test {\ifblank{#2}} and test {\ifblank{#3}} and test {\ifblank{#6}}}{\vspace{4pt}}{%
    \vspace{2pt}%
    {\noindent\fontsize{9.5}{13}\selectfont
      \ifblank{#2}{}{{\color{Ink}#2}}%
      \ifboolexpr{test {\ifblank{#2}} or test {\ifblank{#3}}}{}{\linksep}%
      \ifblank{#3}{}{{\color{Ink}\href{tel:#4}{#3}}}%
      \hfill{\color{Muted}#6}\par}%
    \vspace{5pt}}}
% Width of the row at the current gap and size, into \cvlinkswidth.
\newcommand{\cvlinksmeasure}[1]{%
  \settowidth{\cvlinkswidth}{\fontsize{\cvlinksize}{12}\selectfont#1}}
% One step of the fit by recursion, not \loop: font loading at a new size
% may run a \loop of its own and clobber ours.
\newcommand{\cvlinksfit}[1]{%
  \cvlinksshrinkfalse
  \ifdim\cvlinkswidth>\linewidth
    \ifdim\cvlinkgap>4pt
      \addtolength{\cvlinkgap}{-0.5pt}\cvlinksshrinktrue
    \else\ifdim\cvlinksize>7.6pt
      \addtolength{\cvlinksize}{-0.2pt}\cvlinksshrinktrue
    \fi\fi
  \fi
  \ifcvlinksshrink
    \def\cvlinksnext{\cvlinksmeasure{#1}\cvlinksfit{#1}}%
  \else
    \let\cvlinksnext\relax
  \fi
  \cvlinksnext}
\newcommand{\cvlinks}[1]{{%
  \setlength{\cvlinkgap}{7pt}%
  \setlength{\cvlinksize}{9.2pt}%
  \cvlinksmeasure{#1}%
  \cvlinksfit{#1}%
  \fontsize{\cvlinksize}{12}\selectfont#1\par}\vspace{2pt}}

% ----------------------------------------------------------------------------
% Sections and blocks
% ----------------------------------------------------------------------------
% Section: tracked uppercase label followed by a hairline to the margin.
\newcommand{\sectionlabel}[1]{\MakeUppercase{#1}\hspace{10pt}{\color{Rule}\leaders\vrule height 3.5pt depth -2.85pt\hfill}\kern0pt}
\titleformat{\section}{\labelfont\fontsize{9.2}{12}\selectfont\color{Accent}}{}{0pt}{\sectionlabel}
\titlespacing*{\section}{0pt}{10pt}{5.5pt}
\newcommand{\cvsection}[1]{\Needspace{4\baselineskip}\section*{#1}}

\newcommand{\hl}[1]{{\bfseries\color{Highlight}#1}}

\newcommand{\cvsummary}[2]{#1\par\ifblank{#2}{}{\vspace{4.5pt}#2\par}}

% Skills: bold label and the list on the same line, one line per category.
\newenvironment{cvskills}{\par}{\par}
\newcommand{\cvskill}[2]{%
  \par\hangindent=10pt\hangafter=1\noindent
  {\bfseries\color{Ink}#1:}\hspace{4pt}#2\par\vspace{2.6pt}}

% Job: title | company on the left, dates and location on the right,
% a one-line company description underneath; never stranded at a page end.
%   \cvjob{title}{company}{dates}{location}{company description}
\newcommand{\cvjob}[5]{%
  \par\addvspace{8pt}\Needspace{4\baselineskip}%
  {\fontsize{10.6}{14}\selectfont{\bfseries\color{Ink}#1}\hspace{7pt}{\color{Divider}|}\hspace{7pt}{\bfseries\color{Accent}#2}\hfill{\fontsize{8.8}{14}\selectfont\color{Muted}#3\ifblank{#4}{}{\hspace{5pt}·\hspace{5pt}#4}}}\par
  {\fontsize{8.8}{12.4}\selectfont\color{Muted}#5}\par\nopagebreak
  \vspace{3.5pt}}

\newenvironment{cvbullets}{%
  \begin{itemize}[leftmargin=12.5pt,labelsep=6.5pt,itemsep=2pt,topsep=0pt,parsep=0pt,partopsep=0pt,
    label={\color{Accent}\fontsize{8.5}{8.5}\selectfont\textbullet}]}
  {\end{itemize}}

% Education: one line like a job's, degree | school, the dates (if any) at
% the right edge.
\newcommand{\cvedu}[3]{%
  {\fontsize{10.6}{14}\selectfont{\bfseries\color{Ink}#1}\hspace{7pt}{\color{Divider}|}\hspace{7pt}{\bfseries\color{Accent}#2}\hfill{\fontsize{8.8}{14}\selectfont\color{Muted}#3}}\par}

% Languages share one line, joined by a middle dot. A wide gap instead would
% make parsers read the second language as a separate column.
\newcommand{\cvlangsep}{}
\newcommand{\cvlang}[2]{\cvlangsep{\bfseries\color{Ink}#1:}\hspace{4pt}#2%
  \gdef\cvlangsep{\hspace{6pt}{\color{Muted}·}\hspace{6pt}}}

% ============================================================================
\begin{document}
\fontsize{9.5}{13}\selectfont
\color{Body}
\raggedright
\frenchspacing

\input{\ContentFile}

\end{document}
"
% \ContentFile defaults to "content", \FontDir to "fonts/" (fonts/ibm-plex-sans
% of the repo). tailor-tex does all of this in a temporary folder per build.
% An empty job location drops its middle dot instead of
% leaving it dangling, and a summary with an empty second paragraph drops
% the gap before it.
%
% Font: IBM Plex Sans (OFL, fonts/ibm-plex-sans/OFL.txt). Type scale (pt):
% name 29, headline 12.5, job title 10.6, body 9.5 on 13, metadata 8.8,
% section labels 9.2 tracked. Accent violet #6D28D9.
%
% ATS rules this layout keeps (verified with poppler and a drawing-order
% parser): single text column, real text only, no tables, icons, header or
% footer; no ligatures and no hyphenation; a font that maps every glyph to
% plain ASCII (hyphens and parentheses included); tracking at or below
% 0.06 em; skills one category per line; inline items joined by middle dots,
% never by wide gaps; nothing right-aligned alone where a parser would read
% it as a separate column.
% ============================================================================
\documentclass[a4paper,10pt]{article}

\providecommand{\ContentFile}{content}
\providecommand{\FontDir}{fonts/}

% ----------------------------------------------------------------------------
% Page and packages
% ----------------------------------------------------------------------------
\usepackage[a4paper,left=15mm,right=15mm,top=13mm,bottom=14mm]{geometry}
\usepackage{fontspec}
\usepackage[english]{babel}
\usepackage[protrusion=true]{microtype}
\usepackage{enumitem}
\usepackage{titlesec}
\usepackage{needspace}
\usepackage{etoolbox}
\usepackage{xcolor}
\usepackage[hidelinks,unicode]{hyperref}

% ----------------------------------------------------------------------------
% Fonts: plain letters only, so "fi", "fl" and "ff" stay searchable. IBM Plex
% Sans, bold text in SemiBold; it keeps I, l and 1 distinct, so OCR-based
% checkers read it well.
% ----------------------------------------------------------------------------
\def\CvRegular{IBMPlexSans-Regular.ttf}
\def\CvBold{IBMPlexSans-SemiBold.ttf}
\def\CvHeavy{IBMPlexSans-Bold.ttf}
\def\CvNameTrack{-0.5}
\defaultfontfeatures{Ligatures={TeX,NoCommon},Kerning=On}
\setmainfont{\CvRegular}[Path=\FontDir,BoldFont=\CvBold]
\newfontfamily\namefont{\CvHeavy}[Path=\FontDir,LetterSpace=\CvNameTrack]
% Tracking stays below 0.1 em, so parsers still read every label as one word.
\newfontfamily\labelfont{\CvBold}[Path=\FontDir,LetterSpace=6]

% No hyphenation: every keyword stays whole for ATS parsers.
\hyphenpenalty=10000
\exhyphenpenalty=10000
\setlength{\emergencystretch}{3em}

% ----------------------------------------------------------------------------
% Colors: one violet accent, a deep violet-ink for highlights, neutral text
% ----------------------------------------------------------------------------
\providecommand{\AccentHex}{6D28D9}
\definecolor{Accent}{HTML}{\AccentHex}
\colorlet{Highlight}{Accent!32!black}
\colorlet{Rule}{Accent!18!white}
\colorlet{Divider}{Accent!45!white}
\definecolor{Ink}{HTML}{14171C}
\definecolor{Body}{HTML}{2B3038}
\definecolor{Muted}{HTML}{646B76}
\setlength{\parindent}{0pt}
\setlength{\parskip}{0pt}
\raggedbottom
% Never split a bullet so that one of its lines sits alone on a page.
\clubpenalty=10000
\widowpenalty=10000
\displaywidowpenalty=10000

% No header or footer: running text there lands in the middle of a job's
% bullets when ATS parsers extract the page.
\pagestyle{empty}

% ----------------------------------------------------------------------------
% Header in three rows: the name with the title at the right edge on its
% baseline; then location (with time zone and work format) and the phone,
% with the core stack at the right edge; then every link on one line. Each of
% location, phone and stack may be empty; with all three empty the row is gone.
% Links use the accent so they read as clickable.
%   \cvheader{name}{location}{phone as shown}{phone digits}{headline}{stack}
%   \cvlinks{\cvlink{url}{text}\linksep\cvlink{url}{text}...}
% The links never wrap: a row too wide for the line first gets narrower gaps
% (7pt down to 4pt), then a smaller size (9.2pt down to 7.6pt). Only past
% that floor does it wrap, rather than turn unreadable.
% ----------------------------------------------------------------------------
\newlength{\cvlinkgap}
\setlength{\cvlinkgap}{7pt}
\newlength{\cvlinksize}
\newlength{\cvlinkswidth}
\newif\ifcvlinksshrink
\newcommand{\linksep}{\hspace{\cvlinkgap}{\color{Divider}·}\hspace{\cvlinkgap}}
\newcommand{\cvlink}[2]{\href{#1}{\color{Accent}#2}}
\newcommand{\cvheader}[6]{%
  \hypersetup{pdfauthor={#1}}%
  \noindent{\namefont\fontsize{29}{33}\selectfont\color{Ink}#1}\hfill
  {\fontsize{12.5}{16}\selectfont\bfseries\color{Accent}#5}\par
  \ifboolexpr{test {\ifblank{#2}} and test {\ifblank{#3}} and test {\ifblank{#6}}}{\vspace{4pt}}{%
    \vspace{2pt}%
    {\noindent\fontsize{9.5}{13}\selectfont
      \ifblank{#2}{}{{\color{Ink}#2}}%
      \ifboolexpr{test {\ifblank{#2}} or test {\ifblank{#3}}}{}{\linksep}%
      \ifblank{#3}{}{{\color{Ink}\href{tel:#4}{#3}}}%
      \hfill{\color{Muted}#6}\par}%
    \vspace{5pt}}}
% Width of the row at the current gap and size, into \cvlinkswidth.
\newcommand{\cvlinksmeasure}[1]{%
  \settowidth{\cvlinkswidth}{\fontsize{\cvlinksize}{12}\selectfont#1}}
% One step of the fit by recursion, not \loop: font loading at a new size
% may run a \loop of its own and clobber ours.
\newcommand{\cvlinksfit}[1]{%
  \cvlinksshrinkfalse
  \ifdim\cvlinkswidth>\linewidth
    \ifdim\cvlinkgap>4pt
      \addtolength{\cvlinkgap}{-0.5pt}\cvlinksshrinktrue
    \else\ifdim\cvlinksize>7.6pt
      \addtolength{\cvlinksize}{-0.2pt}\cvlinksshrinktrue
    \fi\fi
  \fi
  \ifcvlinksshrink
    \def\cvlinksnext{\cvlinksmeasure{#1}\cvlinksfit{#1}}%
  \else
    \let\cvlinksnext\relax
  \fi
  \cvlinksnext}
\newcommand{\cvlinks}[1]{{%
  \setlength{\cvlinkgap}{7pt}%
  \setlength{\cvlinksize}{9.2pt}%
  \cvlinksmeasure{#1}%
  \cvlinksfit{#1}%
  \fontsize{\cvlinksize}{12}\selectfont#1\par}\vspace{2pt}}

% ----------------------------------------------------------------------------
% Sections and blocks
% ----------------------------------------------------------------------------
% Section: tracked uppercase label followed by a hairline to the margin.
\newcommand{\sectionlabel}[1]{\MakeUppercase{#1}\hspace{10pt}{\color{Rule}\leaders\vrule height 3.5pt depth -2.85pt\hfill}\kern0pt}
\titleformat{\section}{\labelfont\fontsize{9.2}{12}\selectfont\color{Accent}}{}{0pt}{\sectionlabel}
\titlespacing*{\section}{0pt}{10pt}{5.5pt}
\newcommand{\cvsection}[1]{\Needspace{4\baselineskip}\section*{#1}}

\newcommand{\hl}[1]{{\bfseries\color{Highlight}#1}}

\newcommand{\cvsummary}[2]{#1\par\ifblank{#2}{}{\vspace{4.5pt}#2\par}}

% Skills: bold label and the list on the same line, one line per category.
\newenvironment{cvskills}{\par}{\par}
\newcommand{\cvskill}[2]{%
  \par\hangindent=10pt\hangafter=1\noindent
  {\bfseries\color{Ink}#1:}\hspace{4pt}#2\par\vspace{2.6pt}}

% Job: title | company on the left, dates and location on the right,
% a one-line company description underneath; never stranded at a page end.
%   \cvjob{title}{company}{dates}{location}{company description}
\newcommand{\cvjob}[5]{%
  \par\addvspace{8pt}\Needspace{4\baselineskip}%
  {\fontsize{10.6}{14}\selectfont{\bfseries\color{Ink}#1}\hspace{7pt}{\color{Divider}|}\hspace{7pt}{\bfseries\color{Accent}#2}\hfill{\fontsize{8.8}{14}\selectfont\color{Muted}#3\ifblank{#4}{}{\hspace{5pt}·\hspace{5pt}#4}}}\par
  {\fontsize{8.8}{12.4}\selectfont\color{Muted}#5}\par\nopagebreak
  \vspace{3.5pt}}

\newenvironment{cvbullets}{%
  \begin{itemize}[leftmargin=12.5pt,labelsep=6.5pt,itemsep=2pt,topsep=0pt,parsep=0pt,partopsep=0pt,
    label={\color{Accent}\fontsize{8.5}{8.5}\selectfont\textbullet}]}
  {\end{itemize}}

% Education: one line like a job's, degree | school, the dates (if any) at
% the right edge.
\newcommand{\cvedu}[3]{%
  {\fontsize{10.6}{14}\selectfont{\bfseries\color{Ink}#1}\hspace{7pt}{\color{Divider}|}\hspace{7pt}{\bfseries\color{Accent}#2}\hfill{\fontsize{8.8}{14}\selectfont\color{Muted}#3}}\par}

% Languages share one line, joined by a middle dot. A wide gap instead would
% make parsers read the second language as a separate column.
\newcommand{\cvlangsep}{}
\newcommand{\cvlang}[2]{\cvlangsep{\bfseries\color{Ink}#1:}\hspace{4pt}#2%
  \gdef\cvlangsep{\hspace{6pt}{\color{Muted}·}\hspace{6pt}}}

% ============================================================================
\begin{document}
\fontsize{9.5}{13}\selectfont
\color{Body}
\raggedright
\frenchspacing

\input{\ContentFile}

\end{document}
